\section{Gramática de referencia}\label{sec:gramatica}\label{sec:extension}

En esta sección reúno las reglas vigentes. Los nombres españoles son glosas y los
códigos K son identificadores de catálogo. Ninguno asigna sonidos a la escritura.
LaTeX representa las expresiones; no comprueba automáticamente sus tipos ni su verdad.

\subsection{Tipos y unidades}
\begin{longtable}{@{}p{29mm}>{\raggedright\arraybackslash}p{\dimexpr\linewidth-29mm-2\tabcolsep\relax}@{}}
\toprule Tipo & Uso \\ \midrule\endhead
Término & Objeto referido, variable, dato convencional, concepto, propiedad,
conjunto o lista. Sus subtipos no son intercambiables. \\[4pt]
Proposición & Relación aplicada a sus argumentos, fórmula convencional o
combinación lógica bien formada. Puede ser verdadera o falsa en un contexto. \\[4pt]
Programa & Asignación, secuencia o mientras; describe cambios de estado. \\[4pt]
Unidad superior & Proposición afirmada, término mostrado, programa escrito,
definición, atribución o pregunta. Cada unidad ocupa un bloque separado. \\[4pt]
\bottomrule
\end{longtable}
K01--K08 solos nombran categorías; no identifican un objeto concreto ni afirman
una proposición. Un descriptor de propiedad necesita aplicarse a un objeto
mediante R8 para formar una proposición. Una lista no sirve como conjunto sin
una conversión explícitamente declarada.

\subsection{Contexto, índices y referencias}
Antes de escribir, declaro el dominio, los referentes y las convenciones que
afecten al significado: instante, unidades, criterios de propiedades, marco
espacial o modelo. Los catálogos de relaciones y vocabulario fijan los papeles
de cada argumento.

K10 con índice remite a un objeto de la leyenda. Repetir el índice conserva el
referente. Dos índices distintos no garantizan objetos distintos: esa diferencia
debe declararse si importa. Un índice sin declaración deja el texto incompleto.
El signo distingue los espacios de nombres: referencia 1, relación R1,
concepto C1, propiedad P1 y modelo modal 1 no son el mismo identificador.

\textbf{La leyenda rige todo el documento identificado, también al cambiar de
página.} Cada ficha T constituye un ejemplo independiente con su propio contexto.
Una nueva leyenda local debe indicar dónde comienza y termina; no se permite
reasignar silenciosamente un referente dentro del mismo texto. La falta de
información no autoriza inferir una negación.

\subsection{Grupos, orden y continuidad del texto}
Los operadores se escriben delante de sus argumentos. Se leen desde el exterior
hacia el interior de los grupos. El rectángulo continuo delimita alcance; al
escribir a mano puede sustituirse por corchetes completos. Un encierro por sí
solo no niega: la negación requiere K11. Los contornos internos de los glifos
no son grupos sintácticos.

La barra vertical separa argumentos dentro de un grupo. Cada constructor fija
cuántos recibe y en qué orden. R2, por ejemplo, recibe agente, dirección y dato:
\KalaMostrar{\KalaRelacion{2}{\KalaReferencia{1}\KalaSeparador\KalaReferencia{2}\KalaSeparador\KalaReferencia{3}}}
La leyenda de este ejemplo debe identificar esos tres objetos. Invertirlos
cambia la afirmación o produce argumentos del tipo incorrecto.

Se escribe una unidad superior completa por bloque y se leen los bloques de
arriba abajo. Un salto de página conserva el contexto, pero no conecta dos
programas en una secuencia: esa conexión requiere K21. Conservo cada expresión
con sus grupos cerrados en el mismo bloque; si resulta demasiado larga, uso una
definición previa o una página más ancha. No se introduce puntuación implícita
ni se deja que la cercanía entre dos signos invente una composición.

\subsection{Notación convencional}
El borde discontinuo delimita matemática u otra notación convencional declarada;
el fondo y el color son ayudas sin significado. En su interior rigen las
convenciones del dominio: números, variables, unidades, ecuaciones o fórmulas
químicas, por ejemplo. Una región puede contener un término o una proposición;
el contenido determina su tipo. No se presupone que una letra conserve su
significado entre ejemplos independientes. Los símbolos adoptados, como
\(\in\), siguen siendo convencionales.

\subsection{Conectivos y cuantificadores}
\begin{longtable}{@{}p{31mm}>{\raggedright\arraybackslash}p{\dimexpr\linewidth-31mm-2\tabcolsep\relax}@{}}
\toprule Signo & Argumentos y lectura \\ \midrule\endhead
K11 Negación & Una proposición: niega todo su grupo. \\[4pt]
K12 Conjunción & Dos proposiciones: ambas verdaderas, sin orden temporal. \\[4pt]
K13 Condicional & Antecedente y consecuente: implicación material clásica;
solo es falsa con antecedente verdadero y consecuente falso. \\[4pt]
K17 Alternativa & Dos proposiciones: disyunción inclusiva, verdadera si
al menos una lo es, incluidas ambas. \\[4pt]
K14 Universal & Variable, conjunto dominio y cuerpo proposicional:
el cuerpo se cumple para cada elemento. \\[4pt]
K15 Existencia & Variable, conjunto dominio y cuerpo proposicional:
el cuerpo se cumple para algún elemento. \\[4pt]
\bottomrule
\end{longtable}
Una variable queda vinculada por el cuantificador más cercano que la declare,
solo dentro de su cuerpo, no en el dominio de ese mismo cuantificador.
Evito reutilizar el mismo nombre en cuantificadores anidados. Renombrar variables
vinculadas conserva el significado si no captura otras variables. Los parámetros
libres deben declararse en el contexto. Sobre un dominio vacío, el universal
es verdadero y el existencial falso.

\KalaMostrar{\KalaNegacion{\KalaConjuncion{\KalaMat{P}}{\KalaMat{Q}}}}
\KalaMostrar{\KalaConjuncion{\KalaNegacion{\KalaMat{P}}}{\KalaMat{Q}}}
La primera expresión niega que ambas proposiciones sean verdaderas; la segunda
afirma que P es falsa y Q verdadera. Son T02 y T03. Asimismo, T04 permite un
libro distinto para cada persona; T05 exige un mismo libro para todas. El orden
de cuantificadores forma parte del significado.

La alternativa exclusiva se compone como
\((P\lor Q)\land\neg(P\land Q)\). La negación no crea antónimos y K13 no
expresa causa ni bifurcación de un programa.

\subsection{Colecciones, selección y definiciones}
K06 con tipo 1 enumera un conjunto: el orden y las repeticiones no cambian sus
miembros. Con tipo 2 enumera una lista: ambos pueden importar. El grupo vacío
representa la colección vacía del tipo indicado. R5 requiere un conjunto como
segundo argumento.

K20 recibe variable, conjunto dominio y condición proposicional, y produce
\(\{x\in D\mid\varphi(x)\}\). Vincula la variable solo en la condición.
Construir ese conjunto no afirma que contenga elementos. El capítulo de
\hyperref[sec:composicion]{composición} desarrolla ejemplos y restricciones.

K09 introduce un nombre local y su cuerpo definitorio. Declaro su tipo,
explícitamente en la leyenda o mediante una construcción inequívoca: una
selección, por ejemplo, define un conjunto. El nombre se puede usar desde la
unidad siguiente hasta el final del documento o ámbito declarado. No hay
referencias adelantadas, autorreferencia ni ciclos; tampoco se redefine un
nombre dentro del mismo ámbito. Su expansión evita capturar variables.
Una definición no demuestra un teorema ni ejecuta una asignación.

\subsection{Fuentes y preguntas}
K16 recibe una fuente identificada y una proposición. Atribuye el contenido;
no permite concluirlo como hecho ni asigna una confianza numérica. Por ejemplo,
que un registro diga que un sensor está activo no certifica ese estado.

K19 toma una proposición y pregunta si es verdadera. Una respuesta afirmativa
confirma el contenido y una negativa su negación; sin respuesta no se concluye
ninguna de las dos. Las preguntas abiertas de «quién» o «cuánto» no tienen
todavía un constructor propio.

Definición, atribución y pregunta solo aparecen en el nivel superior. No se
niega una pregunta con K11 ni se conjugan dos atribuciones con K12. Sí puede
atribuirse una conjunción completa o una proposición modal con modelo declarado.

\subsection{Posibilidad con un modelo explícito}
K18 lleva el índice de un modelo. El contexto declara un conjunto no vacío
de estados \(W\), el estado actual \(w\in W\), accesibilidad \(A_m\) y
valoración \(V\) de los átomos en cada estado:
\[
w\models\Diamond_m P\quad\Longleftrightarrow\quad
\exists v\in W:\ (w,v)\in A_m\ \text{y}\ v\models P.
\]
Se admite únicamente el fragmento modal proposicional: sus átomos no tienen
variables libres y no se combinan cuantificadores con modalidades en una
misma fórmula. Todas las modalidades de una fórmula usan el mismo modelo;
pueden anidarse. Sin estados accesibles, la posibilidad es falsa. La necesidad
puede expresarse como \(\neg\Diamond_m\neg P\).

El modelo fija el sentido de «posible»; no significa probable ni demostrado.
T13 muestra un modelo finito. Cambiar el modelo puede cambiar la lectura.

\subsection{Programas y estados}
Los programas tienen esta estructura abstracta; los nombres son metanotación:
\[
\alpha ::= x:=e\mid(\alpha;\alpha)\mid\operatorname{mientras}(b,\alpha).
\]
Las expresiones enteras \(e\) usan constantes, variables inicializadas, suma,
resta y multiplicación. El booleano \(b\) combina comparaciones enteras con
negación, conjunción o alternativa. No incluye fuentes, preguntas,
cuantificadores ni modalidades. Los enteros son matemáticos, sin desbordamiento;
no se incluyen división ni efectos externos.

K04 con índice 1 asigna: evalúa el segundo argumento en el estado anterior y
actualiza la variable del primero. Toda variable leída debe estar inicializada.
K21 ejecuta dos programas de izquierda a derecha; el segundo comienza solo si
el primero termina. K22 comprueba la condición antes del cuerpo, lo ejecuta
si es verdadera y vuelve a comprobarla. Si es falsa al entrar, no ejecuta el
cuerpo. Un programa puede no terminar.

La repetición hasta que \(b\) sea verdadera, con comprobación final, se compone
como \(\alpha;\operatorname{mientras}(\neg b,\alpha)\). Ejecuta el cuerpo
al menos una vez. T15 inicializa \(x=0\) y lo incrementa hasta 3: termina en
\(x=3\). Para límite entero \(n\geq0\) realiza \(n\) incrementos y \(n+1\)
comprobaciones. El coste es \(\Theta(n+1)\) bajo coste unitario; un análisis
con enteros de tamaño arbitrario necesita declarar el coste en bits.

\subsection{Ampliar sin introducir ambigüedad}
Una relación nueva necesita identificador, definición, tipos, cantidad y orden
de argumentos, contexto y ejemplos. Un concepto nuevo necesita dominio y
criterio de pertenencia. Una grafía nueva necesita significado, trazado y
comparación con el inventario vigente. Un giro, altura o modificador no adquiere
significado automáticamente.

Este fragmento no es un vocabulario completo de todas las ciencias ni un
verificador de demostraciones. Sus reglas permiten escribir expresiones
acotadas y ampliar catálogos sin sustituir la notación especializada ya útil.
